(1) 对总动量作正交坐标变换：
\(\mathbf P_{\rm cm}=\frac1{\sqrt N}\sum_i\mathbf p_i\)，另有 \(N-1\) 个三维正交相对动量。
则 \(\delta^{(3)}(\sum_i\mathbf p_i)=N^{-3/2}\delta^{(3)}(\mathbf P_{\rm cm})\)，且动能正交可加。记 \(d=3(N-1)\)，动量可用区域为 \(d\) 维半径 \(\sqrt{2mE}\) 的球。
\[
\boxed{\Omega_G=\frac{V^N}{N!h^{3N}}N^{-3/2}
\frac{\pi^{d/2}(2mE)^{d/2}}{\Gamma(d/2+1)} }.
\]
由固定总动量的Dirac \(\delta\) 造成态数具有“每动量体积”的量纲，若要求熵本身的对数无量纲，可乘一个与 \(E,V\) 无关的固定参考动量体积；它只给熵加常数，不影响下述温度与压强。用 \(h^{3N}\) 和 \(N!\) 表示经典半经典态数规范。

(2) 真实等能量微正则平均采用球壳，而非累积球的体积均匀平均。在 \(d\) 维球壳上，对任一单位向量投影 \(X\)，
\(\langle X^2\rangle=(2mE)/d\)。
固定单粒子单坐标 \(p_{i\alpha}\) 对相对动量子空间的投影范数平方为 \(1-1/N\)。因此
\[
\langle p_{i\alpha}\rangle=0,\qquad
\langle p_{i\alpha}^2\rangle=\frac{2mE}{d}\left(1-\frac1N\right)=\frac{2mE}{3N}.
\]
又因 \(\langle|\sum_i\mathbf p_i|^2\rangle=0\) 和粒子交换对称，
\[
\langle\mathbf p_i^2\rangle=\frac{2mE}{N},\quad
\langle\mathbf p_i\cdot\mathbf p_j\rangle=-\frac{2mE}{N(N-1)}\quad(i\ne j),
\quad \langle p_i^2/(2m)\rangle=E/N.
\]
这种负交叉关联并不意味着粒子之间存在相互作用力，而是由整体守恒条件诱发。

(3) 由于 \(\Omega_G\propto V^N E^{d/2}\)，微正则Gibbs体积熵给
\[
\frac1{T_G}=\frac{\partial S_G}{\partial E}=\frac{d k_B}{2E},
\quad E=\frac d2 k_BT_G,\quad C_{V,G}=\frac d2 k_B.
\]
同时
\[
\frac p{T_G}=\left(\frac{\partial S_G}{\partial V}\right)_E
=\frac{Nk_B}{V},\qquad pV=Nk_BT_G.
\]
由 \(dS_G=0\) 且 \(N\) 固定，\(V^N E^{d/2}=\)常数，故
\(EV^{2N/d}=\)常数，以及
\(TV^{2N/d}=\)常数（由于 \(d=3(N-1)\)）。有限体系中去除质心动量的三个自由度导致温度定义中的 \(3N-3\) 与状态方程压力自由度 \(N\) 同时出现，二者不矛盾。

(4) 令 \(\alpha=d/2=3(N-1)/2\)，则
\(\omega(E)=\alpha\Omega_G/E\propto E^{\alpha-1}\)，
\[
T_B=\frac{E}{(\alpha-1)k_B},\quad
C_{V,B}=(\alpha-1)k_B.
\]
\(T_B\ne T_G\) 属有限体系常见的体积熵与表面熵定义差异；在 \(N\to\infty\) 且 \(E/N\) 固定时相对差异 \(O(1/N)\) 消失。球面上单位投影的四阶矩为 \(3(2mE)^2/[d(d+2)]\)，故
\[
\langle p_{i\alpha}^4\rangle=
\frac{3(2mE)^2(1-1/N)^2}{d(d+2)},\qquad
\frac{\langle p_{i\alpha}^4\rangle}{\langle p_{i\alpha}^2\rangle^2}
=\frac{3d}{d+2}\longrightarrow3.
\]
有限维球壳分布有界而非严格高斯，但随着自由度增长，单个分量趋于正态；这与第二问的总动量负关联在大 \(N\) 下相对减弱一致。
